% !TEX program = xelatex
\documentclass[11pt, a4paper]{article}

\usepackage{fontspec}
\setmainfont[
  Path=/usr/share/texmf/fonts/opentype/public/tex-gyre/,
  UprightFont=texgyrepagella-regular.otf,
  BoldFont=texgyrepagella-bold.otf,
  ItalicFont=texgyrepagella-italic.otf,
  BoldItalicFont=texgyrepagella-bolditalic.otf
]{TeX Gyre Pagella}

\usepackage{amsmath, amssymb, amsthm}
\usepackage[margin=1.25in]{geometry}
\usepackage{setspace}
\onehalfspacing
\usepackage{parskip}
\setlength{\parindent}{0pt}
\setlength{\parskip}{0.85em}
\usepackage[hidelinks]{hyperref}
\usepackage{titlesec}
\titleformat{\section}{\large\bfseries}{}{0em}{}[\vspace{0.2em}\hrule\vspace{0.4em}]
\titleformat{\subsection}{\normalsize\bfseries\itshape}{}{0em}{}
\usepackage{epigraph}
\setlength{\epigraphwidth}{0.65\textwidth}
\renewcommand{\epigraphflush}{center}
\usepackage{booktabs}
\usepackage{array}

\theoremstyle{definition}
\newtheorem{definition}{Definition}
\newtheorem{theorem}{Theorem}
\newtheorem{proposition}{Proposition}
\newtheorem{corollary}{Corollary}
\newtheorem{conjecture}{Conjecture}
\theoremstyle{remark}
\newtheorem{remark}{Remark}

\usepackage{titling}
\setlength{\droptitle}{-2em}

\title{\textbf{From Synchrony to Structure}\\
{\large Coordination Geometry as a General Principle\\
of Adaptive Computation}}
\author{Flyxion \\ \small Independent Researcher}
\date{2026}

\begin{document}

\maketitle

\begin{abstract}
Three recent developments in unconventional computing, taken together,
suggest a common computational principle. In neuromorphic systems,
oscillatory synchrony between network nodes enables concept drift detection
through changes in coordination patterns before those changes are visible
in measurement statistics. In fluid antenna systems, physical communication
structures are continuously reorganized by generative reinforcement learning,
with port selection and beamforming weights jointly optimized to match
changing propagation geometry. In analog compute-in-memory hardware,
block-sparse random routing with input-balanced sign encoding achieves
near-dense performance while reducing inter-core traffic by over~97\%,
with optimal effective fan-in remaining approximately constant as network
scale grows. This essay argues that all three systems may be fruitfully
understood as instances of the same principle: adaptive intelligence
emerges from the continual reorganization of relational geometry rather
than the manipulation of fixed representations. We formalize this perspective
through the Coordination Geometry Framework, in which a system's adaptive
capacity is measured by its ability to reorganize $\mathcal{G}(X)$~--- the
structure of relationships it maintains over its state space~$X$~--- in
response to environmental change. Four claims motivate the analysis:
coordination-first detection, representation as coordination residue,
sparse coordination sufficiency, and small-world scaling. The RSVP field
triple $(\Phi, \mathbf{v}, S_{\mathrm{RSVP}})$ provides a unified
field-theoretic vocabulary with per-input synchrony identified with the
repair ratio~$\rho$ and adaptation with lamphrodyne relaxation. The essay
closes with a conjecture: computation is not fundamentally the manipulation
of symbols or parameters but the continual reorganization of coordination
geometry under constraint.
\end{abstract}

\newpage

\epigraph{\textit{Computation is not the manipulation of symbols under rules.
Computation is the continual reorganization of coordination geometry under
constraint.}}{}

\section{Three Systems, One Candidate Principle}

Consider three recently proposed adaptive systems that appear, at first
reading, to have nothing to do with each other.

The first is a neuromorphic learning architecture in which connection weights
are replaced by oscillatory links whose coupling strengths vary rhythmically
\cite{whitehouse2026}. Learning occurs by adjusting the phases of oscillators
rather than static weights. When the statistical structure of the input
changes, the synchronization pattern changes as well, revealing distributional
shifts before conventional anomaly detectors notice them. The system detects
concept drift not by comparing distributions but by monitoring how its
internal coordination reorganizes around the inputs.

The second is a fluid antenna system in which the physical geometry of a
communication array is continuously reorganized by a generative-adversarial
and reinforcement-learned controller \cite{zhang2026}. Instead of fixed
antenna elements, the system dynamically selects which ports to activate and
what weights to apply, jointly optimizing both decisions to match the
changing propagation environment.

The third is an analog compute-in-memory architecture called ScRRAMBLe
\cite{jaltare2026} in which large neural networks are implemented on
crossbar arrays of non-volatile memory by imposing block-sparse random
routing between computational cores. Rather than dense connectivity,
ScRRAMBLe uses sparse inter-core routing tensors with an input-balancing
constraint that encodes signed weights using only positive conductances.
Networks with as little as 10--25\% connection density match the performance
of fully-connected counterparts while reducing inter-core traffic by over
97\%. Crucially, the effective fan-in --- the average number of incoming
connections per core --- remains approximately constant as the number of
cores grows, a signature of small-world organization.

In the first system, computation proceeds through phase relationships among
coupled oscillators \cite{kuramoto1984, strogatz2003, acebron2005}. In the
second, through dynamic selection among physical configurations. In the
third, through block-sparse routing geometry among memory cores. The three
systems operate in entirely different domains, yet they share a structural
pattern suggestive of a common principle: the relevant computational object
is not a state but a geometry of coordination, and adaptation is not
parameter adjustment but geometric reorganization. This essay attempts to
make that pattern precise, develop its consequences, and assess how far
the parallel extends.

\section{The Legacy of Fixed Computation}

The dominant paradigm in computation rests on an assumption almost never
made explicit: that the computational substrate is fixed while the state it
stores varies. A circuit computes through fixed gates. A neural network
computes through fixed weights. A communication system transmits through
fixed antenna elements. A memory accelerator stores through fixed crossbar
connections. The architecture is invariant; the content is variable.

This assumption makes computation tractable, but reveals its limitation in
changing environments. When the data distribution shifts, the propagation
channel changes, or the problem geometry evolves, a fixed substrate
optimized for prior conditions may fail. The standard response is retraining:
detect the change, discard the old substrate, learn a new one. This is the
cycle of fix, break, refix.

Cybernetics already identified this limitation: a system maintaining stable
behavior in a variable environment must have sufficient internal variety to
absorb the variety of the environment \cite{ashby1956}. If the substrate
can reorganize continuously, adaptation becomes an ongoing process rather
than a periodic crisis.

\begin{definition}[Fixed versus reconfigurable computation]
A \textit{fixed computation} is a mapping $f_W: X \to Y$ parameterized by
a static structure $W$ (weights, topology, physical configuration). Adaptation
requires changing $W$, treated as an exceptional event.

A \textit{reconfigurable computation} is a mapping $f_{\mathcal{G}(t)}: X
\to Y$ parameterized by a dynamically evolving coordination geometry
$\mathcal{G}(t)$. Adaptation is continuous: $\mathcal{G}(t)$ reorganizes
in response to the input stream, and the mapping~$f$ changes accordingly.
\end{definition}

The distinction is not merely the presence of parameter updates. Gradient
descent updates $W$ at each step but is still fixed computation in the
relevant sense because the architecture does not reorganize its topology.
Reconfigurable computation means that the structure of relationships among
computational elements changes: which elements are coupled, how strongly,
and in what configuration.

\section{Rhythmic Sharing: Computation Through Phase}

The Rhythmic Sharing architecture replaces static connection weights with
oscillatory links \cite{kang2026}, inspired by astrocytic oscillatory
dynamics \cite{oneill2023}. The phase matrix $\Phi$ evolves according to:
\[
\frac{d\Phi}{dt} = \omega_0 + (\epsilon_1 + \epsilon_2 \hat{Q}^T \mathbf{n})
\circ \sin(\Psi - \Phi + \gamma),
\]
where $\Psi$ is the local mean phase field, $\hat{Q}$ is the incidence
matrix, $\epsilon_1, \epsilon_2$ are coupling hyperparameters, $\gamma$ is
a phase bias, and $\circ$ denotes element-wise multiplication. The system
uses an Echo State Network substrate \cite{jaeger2004} in which recurrent
dynamics provide a rich manifold of transient states.

The global Kuramoto order parameter
\[
R(t)\,e^{i\langle\Phi\rangle(t)} = \frac{1}{N_l}\sum_{k=1}^{N_l} e^{i\Phi_k(t)}
\]
measures phase coordination across all links \cite{acebron2005}. Per-input
synchrony refines this by measuring coordination locally around each input
channel~$i$, acting as a microscope for channel-specific drift.

The empirical finding is that per-input synchrony is a sensitive early
indicator of distributional change: synchrony for a drifting channel drops
before value-level statistics cross the conventional anomaly threshold.
The coordination is failing before the failure is visible in the measurements.

\begin{proposition}[Coordination-First Detection]
Let $X(t)$ denote an environmental signal and $\mathcal{G}(t)$ the
coordination geometry maintained by an adaptive system. If environmental
change alters the relational structure of $X$ before it alters its marginal
statistics, then perturbations in $\mathcal{G}(t)$ will precede perturbations
in detectors operating directly on $X(t)$. Consequently, monitoring
coordination geometry provides an earlier indicator of system-environment
mismatch than monitoring measurements alone.
\end{proposition}

This proposition is instantiated, but not proved in full generality, by the
Rhythmic Sharing results. It functions here as a motivating hypothesis.

\subsection*{Synchrony as an entropy-reducing coordinate}

The per-input synchrony variable may be interpreted as a local repair ratio.
Let $s_i(t) \in [0,1]$ denote synchrony around input channel~$i$. Define:
\[
S_i(t) = -\log(s_i(t) + \varepsilon),
\]
where $\varepsilon > 0$ prevents singularity at zero synchrony. If a drift
event reduces synchrony from $s_i$ to $s_i - \delta$ with $0 < \delta < s_i$,
then the entropy change is:
\[
\Delta S_i = \log\!\left(\frac{s_i + \varepsilon}{s_i - \delta + \varepsilon}\right) > 0.
\]
A synchrony drop is therefore formally equivalent to an entropy-accessibility
increase. In RSVP terms, this is precisely the transition toward coordination
failure. The detector does not observe the underlying drift directly; it
observes the induced increase in~$S_i$.

\section{Fluid Antennas: Computation Through Reconfiguration}

The fluid antenna system abandons the fixed-position constraint of
conventional arrays \cite{wong2020, wong2024}. The active geometry at
time~$t$ is:
\[
\mathcal{P}(t) = (A(t),\, W(t)),
\]
where $A(t) \in \{0,1\}^{N_s}$ is the binary port activation vector and
$W(t) \in \mathbb{C}^{N_s}$ is the beamforming weight vector. The adaptive
objective is:
\[
\mathcal{P}^*(t) = \arg\max_{\mathcal{P}}\; P(\mathcal{P},\, H(t)),
\]
where $H(t)$ is the propagation channel and $P$ is the performance
functional. The joint optimization is a high-dimensional mixed-integer
non-convex problem solved via a generative reinforcement learning framework
(PD-GPPO) combining a primary-dual network for beamforming with generative
proximal policy optimization for port selection \cite{zhang2026}.

The reorganization dynamics operate at two coupled timescales: beamforming
weights $W(t)$ reconfigure rapidly given a fixed port selection, while port
activation $A(t)$ reconfigures on the episode timescale. Together they
constitute a continuously renegotiated coordination geometry between
transmitter and propagation channel.

\section{ScRRAMBLe: Coordination Through Sparse Routing}

The ScRRAMBLe architecture \cite{jaltare2026} reveals a third instance of
the coordination geometry principle, this time at the hardware level.
Standard analog compute-in-memory accelerators face a scaling problem:
dense connectivity between cores scales quadratically in communication cost.
ScRRAMBLe resolves this by imposing a block-sparse routing structure.

Each core divides its neurons into \textit{slots}~--- subpopulations of
fixed size. Cores communicate through a routing tensor $C_{ijkm}$ that
specifies sparse slot-to-slot connections sampled at a connection density
$p \ll 1$. A critical insight is the input-balancing constraint: for every
positively-weighted incoming connection, a negatively-weighted copy from
the same source is routed to another slot. This ensures that the input to
each core is zero-mean:
\[
\sum_j x_j^{(c)} = 0,
\]
which allows signed weight computation using only positive conductances.
Crucially, the sign is not a property of any individual conductance element;
it is a property of the routing geometry.

\begin{proposition}[Sign from coordination]
\label{prop:sign}
Under the input-balancing constraint, the computation
$y = W^{(c)} x^{(c)}$
with signed weights $W^{(c)}$ is implemented by the positive-conductance
matrix $G^{(c)}$ via the linear relation $y = a \cdot G^{(c)} x^{(c)}$,
where $a$ is a scalar and the zero-mean condition
$\sum_j x_j^{(c)} = 0$
eliminates the bias introduced by the conductance offset.
Consequently, function (signed computation) is not a property of
individual elements but of coordination pattern: the same positive
conductance participates in positive or negative effective weight depending
on its position in the routing geometry.
\end{proposition}

Proposition~\ref{prop:sign} is proved directly by the input-balancing
algebra of \cite{jaltare2026}. Its broader implication is that even the
most elementary computational property~--- sign~--- can be a relational
property of coordination geometry rather than an intrinsic property of a
component.

\subsection*{Small-world coordination}

ScRRAMBLe reveals an unexpected regularity. As network size increases,
the optimal connection density $p^*$ decreases, yet the effective fan-in
$\eta = p^* N_{\mathrm{cores}}$ remains approximately constant. Networks
with 20 cores use $p^* \approx 0.5$; networks with 60 cores use
$p^* \approx 0.2$; in both cases $\eta \approx 10$--12.

This suggests that adaptive systems maintain a characteristic coordination
radius independent of total scale. New cores do not require dense
connectivity to all existing cores; coordination remains localized while
still permitting global information propagation~--- the hallmark of
small-world networks \cite{watts1998, bassett2006}. In the coordination
geometry framework, adaptive intelligence appears to operate near a regime
of constrained connectivity rather than maximal connectivity: intelligence
emerges not from connecting everything to everything else but from
maintaining sufficient pathways for repair and reorganization while
avoiding unnecessary coordination burden.

\subsection*{Routing as CLIO projection}

The ScRRAMBLe routing tensor $C_{ijkm}$ functions as a physically
realizable CLIO projection operator. The full state space of possible
core activations is projected onto a lower-dimensional routing manifold:
\[
\pi_C: X \to \mathcal{G}_{\mathrm{sparse}},
\]
where $\mathcal{G}_{\mathrm{sparse}}$ is the subspace of coordination
geometries compatible with the routing constraints. The environment is not
represented directly; only those relations that survive the routing tensor
are preserved. This makes CLIO no longer merely a philosophical projection
operator but a physically realizable routing architecture instantiated in
hardware.

\section{The Coordination Geometry Framework}

\begin{definition}[Coordination geometry framework]
\label{def:cgf}
A \textit{coordination geometry framework} consists of:
\begin{itemize}
\item A state space $X$ representing the environment the system operates in.
\item A geometry space $\mathcal{G}$ whose elements are coordination
structures~--- configurations of relationships among the system's internal
components.
\item A performance functional $P: \mathcal{G} \times X \to \mathbb{R}$
measuring how well coordination geometry $G \in \mathcal{G}$ serves the
system's objectives given environmental state $x \in X$.
\item A reorganization dynamics $\dot{G} = \mathcal{R}(G, x, \nabla_G P)$
governing how the coordination geometry evolves in response to the
environment and the performance gradient.
\end{itemize}
A system is \textit{adaptively intelligent} to the degree that its
reorganization dynamics $\mathcal{R}$ maintain $G(t)$ in a high-performance
region of $\mathcal{G}$ as $x(t)$ changes.
\end{definition}

\bigskip
\noindent
\begin{tabular}{@{}p{0.20\textwidth}p{0.24\textwidth}p{0.24\textwidth}p{0.24\textwidth}@{}}
\toprule
\textbf{Element} & \textbf{Rhythmic Sharing} & \textbf{Fluid Antenna}
  & \textbf{ScRRAMBLe} \\
\midrule
State space $X$ & Sensor stream & Channel $H$ & Input activations \\
Geometry $\mathcal{G}$ & Phase $\{\Phi,\kappa_{ij}\}$ & Port config $(A,W)$
  & Routing tensor $C_{ijkm}$ \\
Performance $P$ & Synchrony / F1 & Capacity / SNR & Task accuracy \\
Reorganization $\mathcal{R}$ & Phase dynamics & PD-Net + GPPO
  & Training + sparsity \\
Coordination object & Per-input synchrony & Beampattern & Slot routing \\
\bottomrule
\end{tabular}
\bigskip

\begin{conjecture}[Coordination geometry as generalization of parameter adjustment]
\label{conj:generalization}
Any adaptive system whose adaptation can be described as parameter adjustment
$\dot{W} = f(W, X, \nabla_W L)$ can be re-expressed in the coordination
geometry framework by identifying $\mathcal{G}$ with the configuration of
relationships among computational elements induced by~$W$. Parameter
adjustment is then a special case in which the topology of relationships is
fixed and only coupling strengths change, while the coordination geometry
framework additionally permits topological reorganization. This embedding
from parameter space into geometry space has not been formally established
here.
\end{conjecture}

\subsection*{A taxonomy of reorganization types}

Not all coordination geometry reorganizations are equivalent. The three
systems examined suggest a preliminary classification along three axes.

\textit{Topological vs.\ parametric reorganization.} Topological
reorganization changes \emph{which} elements are coupled: port activation
in fluid antennas, routing tensor structure in ScRRAMBLe, effective coupling
topology in Rhythmic Sharing via phase locking. Parametric reorganization
changes \emph{how strongly} existing couplings act: beamforming weights,
oscillator coupling strengths. The core claim of the coordination geometry
framework is that topological reorganization is primary and parametric
reorganization is secondary.

\textit{Timescale.} Rhythmic Sharing's phase dynamics operate at the
oscillation period (fast, continuous). Fluid antenna beamforming weights
reconfigure on a channel coherence timescale (fast, continuous). Port
selection reconfigures on the RL episode timescale (slow, discrete).
ScRRAMBLe's routing tensor is fixed at the architecture level and
reorganizes only via retraining (slowest, discrete). A system's adaptive
capacity depends on which timescales are active: topological reorganization
on fast timescales is more powerful but harder to analyze.

\textit{Local vs.\ global reorganization.} ScRRAMBLe's block-sparse
structure ensures that each routing change is local (one slot-to-slot
connection); global performance emerges from many such local changes.
Rhythmic Sharing's phase dynamics are local (each link adjusts based on
its neighbors) but can produce global synchronization transitions. Fluid
antenna port selection is effectively global (the full activation vector
changes). The small-world result suggests that local reorganization with
sparse global connectivity may be the optimal regime for scalable adaptive
systems.

This taxonomy does not change the core argument: all three systems are
better described by coordination geometry reorganization than by parameter
adjustment alone. But it clarifies that the framework is a family of
mechanisms, not a single mechanism, and that different regimes may require
different analysis tools.

\subsection*{A minimal variational form of coordination repair}

The reorganization dynamics can be written as gradient flow on a mismatch
functional. Let:
\[
\mathcal{E}(G,x) = -D(G,x) + \lambda C(G),
\]
where $D(G,x)$ measures coordination quality between geometry~$G$ and
environmental state~$x$, $C(G)$ is a structural cost penalizing excessive
reconfiguration, and $\lambda > 0$ controls the trade-off between adaptation
and stability. A minimal repair dynamics is:
\[
\dot{G} = -\nabla_G \mathcal{E}(G,x) = \nabla_G D(G,x) - \lambda \nabla_G C(G).
\]
The first term is adaptive pressure; the second is coherence-preserving
inertia.

\begin{proposition}[Monotonic repair under gradient flow]
\label{prop:monotone}
If $x$ is held fixed and $G(t)$ evolves by $\dot{G} = -\nabla_G
\mathcal{E}(G,x)$, then $\mathcal{E}(G(t),x)$ is non-increasing along
trajectories.

\begin{proof}
By the chain rule,
\[
\frac{d}{dt}\mathcal{E}(G(t),x)
= \langle \nabla_G \mathcal{E}(G,x),\dot{G} \rangle
= \langle \nabla_G \mathcal{E}(G,x), -\nabla_G \mathcal{E}(G,x) \rangle
= -|\nabla_G \mathcal{E}(G,x)|^2 \leq 0.
\]
The mismatch functional decreases monotonically unless the system is
already at a critical point.
\end{proof}
\end{proposition}

\begin{remark}[Non-stationary environments]
Proposition~\ref{prop:monotone} assumes a fixed environmental state~$x$.
In realistic settings, $x(t)$ and $G(t)$ change concurrently. The total
time derivative of the mismatch functional then includes an environmental
drift term:
\[
\frac{d}{dt}\mathcal{E}(G(t), x(t))
= -|\nabla_G \mathcal{E}|^2
+ \langle \nabla_x \mathcal{E}(G,x),\, \dot{x} \rangle.
\]
The second term can increase $\mathcal{E}$ even as the system descends the
geometry gradient, because the target moves while repair is underway. The
framework does not guarantee convergence in non-stationary environments;
it provides a tracking bound: if $|\nabla_x \mathcal{E} \cdot \dot{x}| \leq
|\nabla_G \mathcal{E}|^2$, then the net mismatch remains non-increasing.
This is the formal expression of the intuition that adaptation must be
faster than environmental change. The coordination-first detection
result (Proposition~1) is valuable precisely because it extends the window
available for repair: by detecting $\Delta S_i > 0$ before performance
degrades, the system gains lead time in which the inequality above is
more easily satisfied.
\end{remark}

\section{Coordination Geometry and Scale}

The coordination geometry framework does not assume that coordination
occurs at a single scale. A coordination geometry may exist among neurons,
among oscillatory subnetworks, among communication ports, among software
modules, or among hardware cores. The framework is scale-relative.

Let $\mathcal{G}_1, \mathcal{G}_2, \ldots, \mathcal{G}_n$ denote coordination
geometries defined at different organizational levels. A higher-level geometry
may emerge from the lower-level geometry through a projection:
\[
\Pi_{i \to i+1} : \mathcal{G}_i \to \mathcal{G}_{i+1},
\]
forming a tower $\mathcal{G}_1 \to \mathcal{G}_2 \to \cdots \to \mathcal{G}_n$.

Coordination detected at one scale need not be visible at another. A local
coordination failure may remain invisible globally; a global coherence
collapse may emerge from the accumulation of many local failures. The
ScRRAMBLe small-world result instantiates this structure concretely:
slot-level coordination (within-core dense connectivity) is distinct from
core-level coordination (sparse inter-core routing), and the two levels have
different optimal densities. The system achieves global performance through
local coordination rather than global dense connectivity.

\section{Coordination Before Representation}

The standard picture of adaptive intelligence places representation at the
center \cite{clark2016}:
\[
\text{Perception} \to \text{Representation} \to \text{Action.}
\]
The coordination geometry framework suggests an alternative, connecting to
enactivist and dynamical-systems traditions \cite{varela1991, beer2000,
gibson1979}:
\[
\text{Coordination} \to \text{Representation} \to \text{Action.}
\]
Stable coordination produces persistent patterns that can be read out as
representations. Representations are not the foundation on which coordination
is computed; they are the residue of coordination that has been sustained.

In Rhythmic Sharing, the ``representation'' of the input distribution is not
a weight matrix but the current coordination geometry: the phase relationships
and coupling strengths the network has settled into through exposure to the
data. In the fluid antenna system, the actively maintained port configuration
and beamforming weights are the system's best current coordination with the
propagation environment. In ScRRAMBLe, the authors explicitly treat the
outputs of individual cores as population-coded feature vectors maintained
by routing structure rather than by dedicated storage. Features persist
through coordination among subpopulations, not through dedicated symbolic
storage locations.

\begin{definition}[Representation as coordination residue]
\label{def:rep}
A \textit{representation} in the coordination geometry framework is a stable
pattern that persists in a system's coordination geometry $\mathcal{G}(t)$
across perturbations of the input. Formally, a representation is an
equivalence class of coordination geometries that produce the same system
behavior under the current class of inputs. Representations are not stored;
they are maintained.
\end{definition}

\begin{proposition}[Representations as coordination invariants]
\label{prop:attractors}
A coordination geometry $G^*$ is a representation in the sense of
Definition~\ref{def:rep} if it belongs to a stable invariant set~$\mathcal{I}$
of the reorganization dynamics under the current environmental class:
\[
G^* \in \mathcal{I}, \quad
\mathcal{R}(G, x, \nabla_G P)\big|_{x \in \mathcal{X}_0} \in T_G\mathcal{I}
\quad \text{for all } G \in \mathcal{I},
\]
where $T_G\mathcal{I}$ is the tangent space of~$\mathcal{I}$ at~$G$, and
the invariant set is Lyapunov stable with basin of attraction
$\mathcal{N}(\mathcal{I})$.
\end{proposition}

\begin{remark}[Fixed points vs.\ limit cycles]
The generalization from fixed points to invariant sets is necessary because
in practice, coordination attractors are often limit cycles or quasi-periodic
orbits rather than point equilibria. Rhythmic Sharing's phase-locked states
are limit cycles: the phases continue to evolve, but their relative
configuration~--- the pattern that constitutes the representation~---
remains stable. Proposition~\ref{prop:attractors} accommodates this by
requiring only that $\mathcal{R}$ maps the invariant set to itself, not
that any fixed point exists within it. A representation is identified with
the behavioral equivalence class of the invariant set, not with a
particular point on it.

The proposition is a definition-as-theorem: it characterizes what we mean
by ``stable representation'' in the dynamical systems language of the
framework, rather than proving from first principles that representations
must be stable invariants. The empirical content comes from the specific
systems: for Rhythmic Sharing, the claim is that the per-input synchrony
values characterizing a stable input regime form a Lyapunov-stable set
under the Kuramoto-like dynamics; for the fluid antenna system, that the
optimal port configuration for a given channel class forms a stable fixed
point of the RL policy; for ScRRAMBLe, that the routing tensor that
achieves near-optimal performance for a given task is stable under small
perturbations of the input data.
\end{remark}

Memory is a coordination geometry that resists perturbation across long
timescales; forgetting is the decay of that geometry when maintenance is
insufficient. A long-term memory is a repair equilibrium sustained by the
neural coordination structures that support it. The same analysis applies to
institutional memory, scientific knowledge, and semantic coherence~--- each
is a coordination geometry that must be actively maintained to persist
\cite{ashby1956, holland1995}. The free-energy principle offers a related
account in which biological systems minimize prediction error rather than
store explicit representations \cite{friston2010}.

Drift is the event in which $G^*$ ceases to be a stable attractor: its
basin of attraction collapses because~$\mathcal{X}_0$ has changed. Detection
of drift is detection of this collapse~--- a synchrony drop, a beampattern
mismatch~--- before the nominal performance metric has yet reflected it.
Adaptation is the emergence of a new attractor~$G^{**}$ suited to the new
environmental class.

\subsection*{Multiple realizations of representation}

The identification of representations with coordination residues does not
imply a one-to-one correspondence between environmental states and
representations. Let $\pi: X \to \mathcal{M}$ be the projection induced
by a coordination geometry. In general, $\pi(x_1) = \pi(x_2)$ may hold for
distinct $x_1 \neq x_2$: the system compresses many environmental
configurations into a single representational state. Representations are
equivalence classes of environmental situations that induce sufficiently
similar coordination patterns.

This many-to-one structure is not a limitation but a requirement for adaptive
behavior. A system producing a distinct representation for every environmental
variation would fail to generalize. Stable behavior requires compression.
ScRRAMBLe makes this explicit: weight sharing across feature dimensions
(the slot mechanism) deliberately maps multiple incoming signals to the same
stored weights, achieving generalization through structured invariance rather
than explicit case enumeration.

\section{RSVP Interpretation}

The RSVP field triple $(\Phi, \mathbf{v}, S_{\mathrm{RSVP}})$ provides a
continuous field-theoretic language for all three systems under the same
description \cite{freeman2000}. The identification is not a proof that RSVP
uniquely determines the behavior of these systems; it demonstrates that the
field-theoretic concepts provide a natural vocabulary for the coordination
geometry framework.

\begin{remark}[Sign convention]
Throughout this section, $S_{\mathrm{RSVP}} = -\log\rho$ where $\rho$ is
the repair ratio~--- the ratio of semantic support to admissibility threshold
at a field node. This is not thermodynamic entropy but an admissibility
measure: $S_{\mathrm{RSVP}} \leq 0$ characterizes states inside the
admissibility manifold ($\rho \geq 1$) and $S_{\mathrm{RSVP}} > 0$
characterizes states outside it ($\rho < 1$). The zero level set is the
admissibility boundary.
\end{remark}

\begin{definition}[Lamphrodyne relaxation]
\textit{Lamphrodyne relaxation} is the process by which a system reorganizes
its coordination geometry to reduce $S_{\mathrm{RSVP}}$ and restore
admissibility following a coordination failure. If $S_{\mathrm{RSVP}}(x) > 0$
at some field node, lamphrodyne relaxation drives $S_{\mathrm{RSVP}}(x)
\to 0$ by reorganizing $\mathcal{G}$ rather than by changing~$x$ directly.
\end{definition}

\begin{definition}[RSVP-coordination correspondence]
Under the identification ($S_{\mathrm{RSVP}} = -\log\rho$, $\Phi = S_{\mathrm{support}}$,
$|\mathbf{v}_{ij}| = \alpha(i,j)$), the three systems map onto RSVP field
dynamics as follows.

\textit{Rhythmic Sharing}: Per-input synchrony $s_i(t)$ is identified with
$\exp(-S_{\mathrm{RSVP}}(x_i))$. High synchrony is low $S_{\mathrm{RSVP}}$,
meaning tight coordination inside the admissibility region. Phase
reorganization when synchrony drops is lamphrodyne relaxation.

\textit{Fluid Antenna}: The active port configuration concentrates $\Phi$:
activating a port at~$r$ increases $\Phi(r)$, concentrating constraint
propagation at that location. Beamforming weights determine the direction
and strength of $\mathbf{v}$. Port selection is the search for the
configuration minimizing $S_{\mathrm{RSVP}}$ across relevant field nodes.

\textit{ScRRAMBLe}: The routing tensor $C_{ijkm}$ determines which field
nodes are connected by $\mathbf{v}$ and with what strength. Sparse routing
is a sparse~$\mathbf{v}$ field; block structure is spatially coherent
constraint propagation. The input-balancing constraint ensures that the
net entropy flux entering each core is zero-mean, maintaining $S_{\mathrm{RSVP}}
\leq 0$ locally by construction.
\end{definition}

\subsection*{Operational proxies for $S_{\mathrm{RSVP}}$}

The RSVP framework is actionable only if $S_{\mathrm{RSVP}}$ can be computed
online before performance degrades. The following domain-specific proxies
are proposed as computable early indicators, distinct from the lagging
performance metrics (F1, SNR, accuracy) that confirm failure after it occurs.

\textit{Rhythmic Sharing}: $S_i(t) = -\log(s_i(t) + \varepsilon)$ where
$s_i(t)$ is per-input synchrony. This is already computed by the algorithm
and has been empirically validated as a leading indicator in the paper's
datasets.

\textit{Fluid antenna}: A natural proxy is the condition number
$\kappa(\hat{H}(t))$ of the effective channel matrix after port projection,
$\hat{H}(t) = H(t) \cdot A(t)$. High condition number indicates that the
current port configuration poorly spans the channel, i.e., the projection
$\pi_{\mathcal{G}}$ is near-singular~--- a coordination failure that will
produce performance degradation before SNR itself drops below threshold.

\textit{ScRRAMBLe}: A natural proxy is the variance of input sums across
cores, $\sigma^2(\sum_j x_j^{(c)})_c$. Under perfect input balancing this
variance is zero. Growing variance signals that the routing geometry is
failing to maintain its balancing constraint~--- the coordination is
degrading even if task accuracy has not yet declined.

These proxies are consistent with the Coordination-First Detection
principle (Proposition~1): each measures a property of the coordination
geometry rather than a property of the output, and each is computable
in advance of performance collapse.

\begin{proposition}[Adaptation as lamphrodyne relaxation]
In all three systems, adaptation has the structure of lamphrodyne relaxation:
\[
S_{\mathrm{RSVP}}(x) > 0 \;\implies\; \text{reorganize } \mathcal{G}
\;\implies\; S_{\mathrm{RSVP}}(x) \leq 0.
\]
In Rhythmic Sharing, the trigger is a synchrony drop; in the fluid antenna
system, a performance decline; in ScRRAMBLe, gradient-based reorganization
of the routing tensor during training. All three are instances of a system
detecting that it has exited its admissibility region and reorganizing to
return.
\end{proposition}

Connecting back to Proposition~\ref{prop:attractors}: the collapse of a
coordination attractor is precisely the event $S_{\mathrm{RSVP}} > 0$.
Lamphrodyne relaxation re-establishes a stable attractor~$G^{**}$ with
$S_{\mathrm{RSVP}} \leq 0$, completing the adaptation cycle. Systems
monitoring $S_{\mathrm{RSVP}}$ directly will therefore detect attractor
collapse earlier than systems waiting for performance to degrade~--- which
is the Rhythmic Sharing result.

\section{CLIO and Projection Geometry}

Both Rhythmic Sharing and the fluid antenna system perform compression
of environmental information into a lower-dimensional coordination manifold:
\[
\pi: X \to \mathcal{M},
\]
where $\mathcal{M}$ is the manifold of admissible coordination states. In
Rhythmic Sharing, the projection maps sensor stream $X(t) \in \mathbb{R}^n$
to the synchrony manifold $S(t) \in [0,1]^n$. In the fluid antenna system,
the projection maps channel $H \in \mathbb{C}^{N_r \times N_t}$ to the
active configuration $\mathcal{P}(t) = (A(t), W(t))$.

ScRRAMBLe generalizes this: the routing tensor $C_{ijkm}$ implements
$\pi_C: X \to \mathcal{G}_{\mathrm{sparse}}$ as a physically realizable
operator. The full high-dimensional input space is projected onto the sparse
subspace of coordination geometries compatible with the routing constraints.
This is not a fixed projection~--- weight sharing across slots means that
the projection adapts as the routing geometry is learned~--- but its
topological structure (which cores connect to which) is fixed by the random
routing tensor and the sparsity parameter~$p$.

\begin{proposition}[Coordination as adaptive CLIO projection]
\label{prop:clio}
In all three systems, the coordination geometry $\mathcal{G}(t)$ functions
as a CLIO projection of the environmental state $X(t)$: a lower-dimensional
manifold that amplifies structure relevant to the system's objectives while
compressing structure that is irrelevant. The projection is adaptive: as
$X(t)$ changes, $\mathcal{G}(t)$ reorganizes to maintain the quality of
the projection. This is the key distinction from fixed dimensionality
reduction: the mapping is continuously renegotiated with the environment,
not applied statically to it.
\end{proposition}

The adaptive CLIO projection is what allows all three systems to detect or
respond to changes that fixed projections would miss. In ScRRAMBLe, the
finding that performance is approximately invariant across a wide range of
connection densities (from 10\% to 80\% in the CIFAR experiments) implies
that many distinct routing geometries project into equivalent coordination
manifolds for the task at hand~--- a form of projection robustness consistent
with Proposition~\ref{prop:clio}.

\section{Coordination, Performance, and Failure}

The coordination geometry framework does not claim that every coordination
geometry is beneficial. Coordination can become maladaptive.

Let $G^*$ be a stable coordination attractor. Stability alone does not imply
optimality. A system may remain trapped in a geometry that was previously
successful but is no longer appropriate: this is the coordination-geometry
analogue of the overfitting problem. A network may develop a stable phase
pattern suited to training data that fails to generalize; a fluid antenna may
lock onto a port configuration suited to a prior channel state; a reservoir
may settle into a dynamics that was once adaptive but no longer matches
the current environment.

The performance functional $P(G, x)$ mediates the trade-off between
stability and reconfigurability. Reorganization dynamics must preserve enough
continuity to maintain useful coordination while remaining flexible enough
to discover new geometries when environmental conditions change.
Proposition~\ref{prop:monotone} shows that gradient flow on the mismatch
functional monotonically reduces~$\mathcal{E}$, but does not guarantee
finding a global optimum. Local attractors in coordination space can
trap the system.

The central problem of adaptive computation is therefore not maximizing
stability or maximizing change. It is maintaining the appropriate balance
between persistence and reorganization~--- between
$\lambda \nabla_G C(G)$ (coherence-preserving inertia) and
$\nabla_G D(G,x)$ (adaptive pressure) in the variational form.

\section{Toward Reconfigurable Intelligence}

The coordination geometry framework, motivated by three concrete systems
from different domains, suggests a broader conjecture about adaptive
computation. Current intelligent systems operate predominantly on fixed
substrates: large language models have fixed weights, base stations have
fixed antenna elements, memory accelerators have fixed crossbar connections,
and when the environment changes in ways the fixed substrate cannot
accommodate, the intelligence fails.

Reconfigurable intelligence operates differently. The substrate reorganizes,
and intelligence is the process of maintaining coordination across that
reorganization. Each of the three systems examined here instantiates a
specific element of Definition~\ref{def:cgf}.

\textit{Rhythmic Sharing in silicon}: Neuromorphic hardware based on
oscillator networks implements Rhythmic Sharing in silicon \cite{maass2002,
lukosevicius2009}. The coordination geometry is the physical oscillator
coupling state ($\mathcal{G}$), and $\mathcal{R}$ operates on phase
relationships rather than stored weights.

\textit{Reconfigurable intelligent surfaces}: These implement the fluid
antenna principle at large scale \cite{direnzo2019, basar2019}.
The surface geometry is $\mathcal{G}$, the propagation channel is~$X$,
and the configuration policy is~$\mathcal{R}$.

\textit{Scalable analog memory}: ScRRAMBLe demonstrates that sparse
coordination geometry can match dense representation at a fraction of
communication cost~--- that the routing manifold $\mathcal{G}_{\mathrm{sparse}}$
is functionally equivalent to the full representation space for the relevant
tasks while requiring 97\% less traffic. This is empirical evidence that
$\mathcal{G}$ can be dramatically sparser than~$X$ without loss of adaptive
capacity.

\textit{Reservoir computing}: A high-dimensional dynamical system projects
inputs onto a rich manifold of transient states, with learning only at the
readout layer \cite{maass2002, jaeger2001, verstraeten2007}. Even when
$\mathcal{G}$ is fixed, coordination geometry provides adaptive capacity
that static mapping cannot; as $\mathcal{G}$ itself becomes adaptive, the
capacity grows further.

\textit{Active inference}: Organisms maintain internal generative models and
minimize free energy by acting on the environment \cite{friston2010}. The
generative model is a coordination geometry and free energy minimization is
lamphrodyne relaxation renegotiating both internal geometry and environmental
coupling simultaneously.

\begin{conjecture}[Reconfigurable Intelligence]
An adaptively intelligent system $\mathcal{N}$ operating in a changing
environment $X(t)$ maintains performance by continuously reorganizing its
coordination geometry $\mathcal{G}(t)$ to satisfy:
\[
S_{\mathrm{RSVP}}(\pi_{\mathcal{G}}(x)) \leq 0 \quad
\text{for all } x \in X(t),
\]
where $\pi_{\mathcal{G}}$ is the CLIO projection induced by the current
coordination geometry and $S_{\mathrm{RSVP}} \leq 0$ is the admissibility
condition. When this condition is violated, the system reorganizes~$\mathcal{G}$
via lamphrodyne relaxation. Intelligence is the capacity to sustain this
reorganization-repair cycle across the range of environmental changes the
system encounters.
\end{conjecture}

\section{Implementation Pathways}

The coordination geometry framework suggests a generic four-layer
architecture for reconfigurable adaptive systems. The common requirement
is that the system must expose its coordination geometry as an object of
measurement and control, rather than hiding it inside fixed parameters.

The \textit{sensing layer} receives an environmental stream $X(t)$~---
sensor data, channel measurements, user behavior, machine telemetry,
biological recordings. Its role is not to solve the task directly but to
preserve enough temporal and relational structure for downstream coordination
dynamics to detect changes that may not yet appear in local statistics.

The \textit{coordination layer} maintains a dynamic geometry $\mathcal{G}(t)$
over the incoming stream: a graph of coupled oscillators, a reservoir state,
an attention topology, a port-selection manifold, a routing tensor, a
reconfigurable sensor layout. The essential requirement is that
$\mathcal{G}(t)$ changes when the relational structure of $X(t)$ changes.

The \textit{reorganization layer} updates $\mathcal{G}(t)$ in response to
mismatch~--- by reinforcement learning, gradient flow, evolutionary search,
message passing, Bayesian filtering, active inference, adversarial generation,
or local Hebbian updates. The specific method is less important than the
fact that the adaptive variable is the system's relational structure, not
merely a parameter vector.

The \textit{readout layer} converts the current coordination geometry into
task-relevant action: a synchrony score, a beamforming pattern, a motor
policy, a routing decision, a stable representation.

A generic implementation:
\[
X(t) \longrightarrow \mathcal{G}(t) \longrightarrow y(t),
\quad\text{with}\quad
\dot{\mathcal{G}}(t) = \mathcal{R}\!\left(\mathcal{G}(t), X(t), P(\mathcal{G}(t),X(t))\right).
\]
The practical design rule: do not only ask what state the system is in. Ask
what coordination geometry the system is maintaining, how that geometry
fails, and what reorganization process restores it.

\section{Open Questions}

The coordination geometry framework currently sits at the level of a research
programme rather than a fully formalized theory. Three questions mark the
boundary between what the framework establishes and what it leaves open.

\textit{Falsifiability.} The claim that adaptive intelligence is
fundamentally coordination geometry reorganization would be falsified by
a system that: (a) achieves robust performance across a changing environment
without any topological reorganization, using only parametric adjustment of
fixed weights; and (b) cannot be re-described as coordination geometry
reorganization without making the framework vacuous (i.e., without
identifying ``geometry'' with the weights themselves). Deep learning systems
that generalize across distribution shifts via weight interpolation are the
strongest candidate counterexample. If interpolation between fixed-weight
solutions is sufficient for robust generalization, the topological
reorganization requirement may be too strong. The framework would need to
be weakened to: ``topological reorganization is \emph{sufficient} but not
\emph{necessary} for robust adaptation, and it is the most efficient mechanism
when it is available.''

\textit{The small-world constant.} The ScRRAMBLe result that effective fan-in
$\eta \approx 10$--12 remains constant across scales is striking. Whether
this constant can be derived from first principles within the coordination
geometry framework is an open problem. A candidate derivation: $\eta$ is the
value at which the repair-rate term $|\nabla_G \mathcal{E}|^2$ in
Proposition~\ref{prop:monotone} balances the environmental drift term
$|\nabla_x \mathcal{E} \cdot \dot{x}|$ at a typical rate of environmental
change. If so, $\eta$ is not a universal constant but a function of the
ratio of reorganization speed to environmental change rate~--- a predicted
scaling law that could be tested by varying the non-stationarity of the
input in ScRRAMBLe-style experiments.

\textit{A common algorithmic principle.} The three systems use different
optimization paradigms: phase-coupled ODEs (Rhythmic Sharing), generative
RL (fluid antennas), gradient descent on routing tensors (ScRRAMBLe). A
candidate unifying principle is online projection pursuit: each system is
continuously searching for the low-dimensional coordination manifold that
best projects the environmental structure relevant to the task.
Rhythmic Sharing finds this manifold through synchronization; fluid antennas
through port selection and beamforming; ScRRAMBLe through sparse routing.
The three are instances of different algorithms for the same optimization
problem: minimize $S_{\mathrm{RSVP}}(\pi_{\mathcal{G}}(x))$ over the
space of admissible coordination geometries. Whether this can be given a
unified algorithmic treatment~--- for instance, as geodesic flow on
$\mathcal{G}$ equipped with a metric derived from the mismatch functional
$\mathcal{E}$~--- is the deepest open question raised by the framework.

\section{Conclusion}

The Rhythmic Sharing algorithm shows that concept drift is detectable through
coordination geometry reorganization before it is detectable in value-level
statistics. The fluid antenna system shows that communication performance
is improvable through physical geometry reorganization beyond what
fixed-substrate optimization achieves. ScRRAMBLe shows that block-sparse
coordination geometry can match dense representation at a fraction of the
communication cost, and that optimal effective fan-in remains constant as
network scale grows~--- an empirical signature of intelligence emerging from
sparse relational structure rather than dense accumulation.

The four central claims of the essay are:

\textit{Coordination-first detection} (Proposition~1): adaptive systems often
detect environmental change through perturbations in coordination geometry
before those perturbations are visible in measurement statistics.

\textit{Representation as coordination invariant} (Definition~3,
Proposition~4): representations are stable invariant sets in coordination
space, accommodating both fixed points and limit cycles. Memory is
coordination maintenance. Forgetting is coordination decay. Drift is
invariant set collapse.

\textit{Sparse coordination sufficiency} (Proposition~\ref{prop:sign},
ScRRAMBLe results): coordination geometries far sparser than the full
representation space can preserve adaptive capacity. Intelligence depends
on the geometry of admissible coordination pathways, not on maximal
connectivity.

\textit{Small-world scaling}: optimal adaptive systems maintain approximately
constant effective fan-in as they scale, operating near a regime of
constrained connectivity rather than maximal connectivity.

The RSVP framework provides a unified vocabulary: $\Phi$ encodes
coordination density, $\mathbf{v}$ encodes its propagation, and
$S_{\mathrm{RSVP}} = -\log\rho$ measures admissibility. Drift is
$S_{\mathrm{RSVP}} > 0$. Adaptation is lamphrodyne relaxation.

Representation is not the foundation of intelligence but its stable residue.
What persists when coordination is sustained is a representation; what
dissolves when coordination fails is a representation. The foundation is
the coordination itself: the ongoing work of maintaining a geometry of
relationships that keeps the system within its admissibility manifold as
the world changes.

The map does not precede the territory. The map is what the territory looks
like when the coordination between observer and environment has been
sustained long enough to leave a stable trace.

\bigskip
\noindent\rule{0.4\textwidth}{0.4pt}

\vspace{0.5em}
\noindent{\small Flyxion / Independent Researcher. 2026.}

\section*{References}

\begin{thebibliography}{99}

\bibitem{whitehouse2026}
Whitehouse, I., Kang, H., and Losert, W.
``Emergent detection of concept drift within the glia-inspired `rhythmic
sharing' algorithm.''
\textit{npj Unconventional Computing} 3:25, 2026.

\bibitem{zhang2026}
Zhang, J., Xing, C., Tang, J., Yang, Z., Deng, N., Zhao, N., Wong, K.-K.,
Niyato, D., and Chae, C.-B.
``Generative reinforcement learning for fluid antenna systems based sensing
and communication.''
\textit{npj Wireless Technology} 2:30, 2026.

\bibitem{jaltare2026}
Jaltare, V., Mondal, R., Gibb, L., Song, J., Leugering, J., and
Cauwenberghs, G.
``ScRRAMBLe: block-sparse deep learning architecture for analog in-memory
computing accelerators.''
\textit{npj Unconventional Computing} 3:26, 2026.

\bibitem{kang2026}
Kang, H., and Losert, W.
``Rhythmic sharing: a bio-inspired paradigm for zero-shot adaptive learning
in neural networks.''
\textit{Physical Review Research} 8:013267, 2026.

\bibitem{oneill2023}
O'Neill, K.~M., et al.
``Decoding natural astrocyte rhythms: dynamic actin waves result from
environmental sensing by primary rodent astrocytes.''
\textit{Advanced Biology} 7:e2200269, 2023.

\bibitem{kuramoto1984}
Kuramoto, Y.
\textit{Chemical Oscillations, Waves, and Turbulence}.
Springer, 1984.

\bibitem{strogatz2003}
Strogatz, S.~H.
\textit{Sync: The Emerging Science of Spontaneous Order}.
Hyperion, 2003.

\bibitem{acebron2005}
Acebr\'{o}n, J.~A., Bonilla, L.~L., P\'{e}rez~Vicente, C.~J.,
Ritort, F., and Spigler, R.
``The Kuramoto model: a simple paradigm for synchronization phenomena.''
\textit{Reviews of Modern Physics} 77(1):137--185, 2005.

\bibitem{jaeger2001}
Jaeger, H.
``The echo state approach to analysing and training recurrent neural
networks.''
Technical Report GMD 148, German National Research Center for Information
Technology, 2001.

\bibitem{jaeger2004}
Jaeger, H., and Haas, H.
``Harnessing nonlinearity: predicting chaotic systems and saving energy in
wireless communication.''
\textit{Science} 304:78--80, 2004.

\bibitem{maass2002}
Maass, W., Natschl\"{a}ger, T., and Markram, H.
``Real-time computing without stable states: a new framework for neural
computation based on perturbations.''
\textit{Neural Computation} 14(11):2531--2560, 2002.

\bibitem{lukosevicius2009}
Luko\v{s}evi\v{c}ius, M., and Jaeger, H.
``Reservoir computing approaches to recurrent neural network training.''
\textit{Computer Science Review} 3(3):127--149, 2009.

\bibitem{verstraeten2007}
Verstraeten, D., Schrauwen, B., D'Haene, M., and Stroobandt, D.
``An experimental unification of reservoir computing methods.''
\textit{Neural Networks} 20(3):391--403, 2007.

\bibitem{wong2020}
Wong, K.-K., Shojaeifard, A., Tong, K.-F., and Zhang, Y.
``Fluid antenna systems.''
\textit{IEEE Transactions on Wireless Communications} 20:1950--1962, 2020.

\bibitem{wong2024}
Wang, C., et al.
``AI-empowered fluid antenna systems: opportunities, challenges, and future
directions.''
\textit{IEEE Wireless Communications} 32:166--173, 2024.

\bibitem{direnzo2019}
Di~Renzo, M., et al.
``Smart radio environments empowered by reconfigurable intelligent surfaces.''
\textit{EURASIP Journal on Wireless Communications and Networking}, 2019.

\bibitem{basar2019}
Ba\c{s}ar, E., Di~Renzo, M., de~Rosny, J., Debbah, M., Alouini, M.-S.,
and Zhang, R.
``Wireless communications through reconfigurable intelligent surfaces.''
\textit{IEEE Access} 7:116753--116773, 2019.

\bibitem{watts1998}
Watts, D.~J., and Strogatz, S.~H.
``Collective dynamics of `small-world' networks.''
\textit{Nature} 393:440--442, 1998.

\bibitem{bassett2006}
Bassett, D.~S., and Bullmore, E.
``Small-world brain networks.''
\textit{Neuroscientist} 12(6):512--523, 2006.

\bibitem{dalgaty2024}
Dalgaty, T., et al.
``Mosaic: in-memory computing and routing for small-world spike-based
neuromorphic systems.''
\textit{Nature Communications} 15:142, 2024.

\bibitem{friston2010}
Friston, K.
``The free-energy principle: a unified brain theory?''
\textit{Nature Reviews Neuroscience} 11(2):127--138, 2010.

\bibitem{freeman2000}
Freeman, W.~J.
\textit{Neurodynamics: An Exploration in Mesoscopic Brain Dynamics}.
Springer, 2000.

\bibitem{beer2000}
Beer, R.~D.
``Dynamical approaches to cognitive science.''
\textit{Trends in Cognitive Sciences} 4(3):91--99, 2000.

\bibitem{clark2016}
Clark, A.
\textit{Surfing Uncertainty: Prediction, Action, and the Embodied Mind}.
Oxford University Press, 2016.

\bibitem{varela1991}
Varela, F.~J., Thompson, E., and Rosch, E.
\textit{The Embodied Mind}.
MIT Press, 1991.

\bibitem{gibson1979}
Gibson, J.~J.
\textit{The Ecological Approach to Visual Perception}.
Houghton Mifflin, 1979.

\bibitem{ashby1956}
Ashby, W.~R.
\textit{An Introduction to Cybernetics}.
Chapman and Hall, 1956.

\bibitem{holland1995}
Holland, J.~H.
\textit{Hidden Order: How Adaptation Builds Complexity}.
Addison-Wesley, 1995.

\end{thebibliography}

\end{document}

